\section{Conclusion}\label{sec:conclusion}
The goal of this thesis is to design and implement a full-stack web framework focused on simplifying the development,
supporting extensibility, and integrating AI-Generation features.
One of the sub goals is to implement the framework with minimal usage of third-party code.
Currently, the prototype does not rely on external libraries or other frameworks, it uses internally developed libraries
and standardized PHP extensions such as curl, PDO, fileinfo, gd, and mbstring.
This fulfills the NFR4.

\subsection*{Hypothesis Evaluation}\label{subsec:hypothesis-evaluation}
The observations during implementation and evaluation results generally support the hypotheses proposed in Chapter~\ref{subsec:introduction-hypothesis}.

The scenarios demonstrated that unified and convention-driven abstractions simplifies the deployment of applications with varying types.
The built-in features enabled rapid deployment of smaller to medium-sized applications while supporting custom implementation if needed.
This and evaluation results qualitatively satisfy the NFR1.
The scenarios confirmed the usability of the tested systems and AI-Generated content.
The System Usability Scale evaluation of the prototype produced an average of 70, which makes the implementation acceptable.

Observations during implementation showed that minimizing reliance on third-party dependencies improved consistency
and streamlined integration of subsystems across the architecture fulfilling the NFR2,
although it increased the implementation time.

The evaluation results show that AI-Generated content decreases the time spent on content creation.
However, the evaluation confirmed the AI-Generated content requires manual verification and modification.

Due to the limited scope of the implemented prototype and smaller number of participants in evaluation,
the conclusion should be interpreted as observation rather than statistically validated result.

\subsection*{Requirements Evaluation}\label{subsec:requirements-evaluation}
The implementation observations and evaluation results indicate that the framework fulfills the majority of requirements.

The NFR3 requirement is perceived as partially satisfied based on qualitative evaluation.
During implementation and evaluation actions performed by the framework did not exceed the expect time to fulfill the request,
thus the framework is performant enough for small-to-medium sized projects.
However, no performance benchmarking was conducted.

\subsection*{Limitations of the Prototype}\label{subsec:limitations-of-the-prototype}
The thesis also identified several limitations of the current implementation.
Currently, some advanced features expected from modern full-stack framework require manual implementation.
These include on-demand user sign in, collaborative editing, e-commerce support, and general visual content editor.
Additionally, the implemented architecture favors solutions maximizing for simplicity, abstraction, and control.
For example, the closed user system simplifies administration but limits the usability in systems opened the public registrations.

\subsection*{AI-Generated Content}\label{subsec:ai-generated-content}
The AI-Generated content demonstrated usefulness for certain tasks, such as article content generation and phrase translation.
Although, the evaluation shows that AI-Generated outputs require manual verification and modification.
As a result, AI-Generation features should currently be considered a supporting tool for drafting purposes rather
than automated solution generation tool.

\subsection*{Future Work}\label{subsec:future-work}
Future Work should focus primarily on improving extensibility of the system, reducing manual implementations, and refactoring
outdated systems.
Potential improvements include replacement of component-based widget editor for general visual content editor,
support for modules, improvement of authentication system, refactoring the structure of the repository, and adding more
page templates such as contact page.
Additionally, conducting a larger-scale usability evaluation with framework-specific questions to identify the weaker parts of the system.

\subsection*{Final Thoughts}\label{subsec:final-thoughts}
Overall, the thesis shows that the implemented framework is usable for wide range of use cases and enables rapid deployment
of simple projects.
Although, the framework currently contains several limitations, the implemented prototype uses solid architectural foundation for future expansions.
